\documentclass[a4paper,12pt]{article}
\usepackage[a4paper,margin=1in]{geometry}
\usepackage{amsmath,amssymb}
\begin{document}
\noindent
\textbf{UJIAN TENGAH SEMESTER (UTS) 2024} \\
\textbf{Mata Kuliah: AS2001 - Astrofisika} \\
\rule{\textwidth}{1pt}

\vspace{1em}
\noindent\textbf{Petunjuk:} Jawablah semua pertanyaan di bawah ini dengan menyertakan penurunan rumus dan substitusi numerik.

\vspace{1.5em}
\noindent\textbf{Soal 1} \\
Sebuah bintang deret utama memiliki radius $R = 2 R_\odot$ dan temperatur efektif permukaan $T_{\text{eff}} = 10000\text{ K}$.
\begin{itemize}
    \item[a.] Tentukan fluks pancaran radiasi total pada permukaan bintang tersebut menggunakan hukum Stefan-Boltzmann.
    \item[b.] Hitung luminositas total bintang tersebut dalam satuan watt (W) dan dalam satuan luminositas matahari ($L_\odot$).
\end{itemize}

\vspace{2em}
\noindent\textbf{Soal 2} \\
Sebuah bintang neutron memiliki massa $M = 1.4 M_\odot$ dan radius rata-rata $R = 10\text{ km}$.
\begin{itemize}
    \item[a.] Tentukan besar kecepatan lepas ($v_{\text{esc}}$) dari permukaan bintang neutron tersebut.
    \item[b.] Nyatakan kecepatan lepas tersebut dalam persentase terhadap kecepatan cahaya ($c$).
\end{itemize}

\end{document}
